% Part: incompleteness
% Chapter: incompleteness-provability
% Section: fixed-point-lemma

\documentclass[../../../include/open-logic-section]{subfiles}

\begin{document}

\olfileid[pt]{inc}{inp}{fix}

\olsection{O lema do ponto fixo}

\begin{explain}
O lema do ponto fixo diz que, para toda !!{formula}~$!B(x)$, existe
!!a{sentence}~$!A$ tal que $\Th{T} \Proves !A \liff !B(\gn{!A})$,
desde que $\Th{T}$ estenda~$\Th{Q}$. No caso da sentença do mentiroso,
queremos que $!A$ seja equivalente (de modo demonstrável em~$\Th{T}$) a
``$\gn{!A}$ é falsa'', isto é, à afirmação de que $\Gn{!A}$ é o número de
G\"odel de uma !!{sentence} falsa. Para compreender a ideia da prova, será
útil compará-la com a glosa informal de Quine para~$!A$: ``{}`produz uma
falsidade quando precedida por sua própria citação' produz uma falsidade
quando precedida por sua própria citação''. A operação de tomar uma
expressão e formar uma sentença precedendo essa expressão por sua própria
citação pode ser chamada de \emph{diagonalização} da expressão, e o
resultado, de sua diagonalização. Assim, a diagonalização de `produz uma
falsidade quando precedida por sua própria citação' é ``{}`produz uma
falsidade quando precedida por sua própria citação' produz uma falsidade
quando precedida por sua própria citação''. Observe agora que a sentença do
mentiroso de Quine não é a diagonalização de `produz uma falsidade', mas de
`produz uma falsidade quando precedida por sua própria citação'. Portanto,
a própria propriedade diagonalizada para produzir a sentença do mentiroso
envolve diagonalização!{}

Na linguagem da aritmética, formamos citações de !!a{formula} com uma
variável livre calculando seu número de G\"odel e substituindo a variável
livre pelo numeral padrão desse número de G\"odel. A diagonalização
de~$!E(x)$ é $!E(\num{n})$, onde $n = \Gn{!E(x)}$. (Daqui em diante,
abreviaremos $\num{\Gn{!E(x)}}$ por $\gn{!E(x)}$.) Assim, se $!B(x)$ é
``é uma falsidade'', então ``produz uma falsidade se precedida por sua
própria citação'' seria ``produz uma falsidade quando aplicada ao número de
G\"odel de sua diagonalização''. Se tivéssemos um símbolo~$\Obj{diag}$ para
a função $\fn{diag}(n)$ que calcula o número de G\"odel da diagonalização da
!!{formula} com número de G\"odel~$n$, poderíamos escrever $!E(x)$ como
$!B(\Obj{diag}(x))$. E a versão de Quine da sentença do mentiroso seria
então sua diagonalização, isto é, $!E(\gn{!E(x)})$ ou
$!B(\Obj{diag}(\gn{!B(\Obj{diag}(x))}))$. Evidentemente, $!B(x)$ poderia
agora ser qualquer outra propriedade, e a mesma construção funcionaria.
Para o teorema da incompletude, tomaremos $!B(x)$ como ``$x$~não é
!!{derivable} em~$\Th{T}$.'' Então $!E(x)$ seria ``produz
!!a{sentence} não !!{derivable} em~$\Th{T}$ quando aplicada ao número de
G\"odel de sua diagonalização''.

Para formalizar isso em~$\Th{T}$, precisamos encontrar uma maneira de
formalizar $\fn{diag}$. A função $\fn{diag}(n)$ é computável; de fato, é
recursiva primitiva: se $n$ é o número de G\"odel de uma fórmula~$!E(x)$,
$\fn{diag}(n)$ devolve o número de G\"odel de~$!E(\gn{!E(x)})$. (Lembre-se:
$\gn{!E(x)}$ é o numeral padrão do número de G\"odel de~$!E(x)$, isto é,
$\num{\Gn{!E(x)}}$.) Se $\Obj{diag}$ fosse um símbolo de função em
$\Th{T}$ que representasse a função $\fn{diag}$, poderíamos tomar $!A$
como a fórmula $!B(\Obj{diag}(\gn{!B(\Obj{diag}(x))}))$. Observe que
\begin{align*}
\fn{diag}(\Gn{!B(\Obj{diag}(x))}) & = 
\Gn{!B(\Obj{diag}(\gn{!B(\Obj{diag}(x))}))} \\
& = \Gn{!A}.
\end{align*}
Supondo que $\Th{T}$ possa !!{derive}
\[
\Obj{diag}(\gn{!B(\Obj{diag}(x))}) = \gn{!A},
\]
ela pode !!{derive} $!B(\Obj{diag}(\gn{!B(\Obj{diag}(x))}))
\liff !B(\gn{!A})$. Mas o lado esquerdo é, por definição,~$!A$.

É claro que $\Obj{diag}$ em geral não será um símbolo de função de
$\Th{T}$ e certamente não é um símbolo de~$\Th{Q}$. Mas, como
$\fn{diag}$ é computável, ela é \emph{representável} em~$\Th{Q}$ por alguma
fórmula $!D_{\fn{diag}}(x,y)$. Assim, em vez de escrever
$!B(\Obj{diag}(x))$, podemos escrever
$\lexists[y][(!D_{\fn{diag}}(x,y) \land !B(y))]$. De resto, a prova
esboçada acima funciona e, na verdade, já funciona em~$\Th{Q}$.
\end{explain}

\begin{lem}
\ollabel{lem:fixed-point} Seja $!B(x)$ qualquer fórmula com uma variável
livre~$x$. Então existe uma sentença~$!A$ tal que $\Th{Q} \Proves
!A \liff !B(\gn{!A})$.
\end{lem}


\begin{proof}
Dada $!B(x)$, seja $!E(x)$ a fórmula
$\lexists[y][(!D_{\fn{diag}}(x,y) \land !B(y))]$ e seja $!A$~sua
diagonalização, isto é, a fórmula $!E(\gn{!E(x)})$.

Como $!D_{\fn{diag}}$ representa $\fn{diag}$, e
$\fn{diag}(\Gn{!E(x)}) = \Gn{!A}$, $\Th{Q}$ pode !!{derive}
\begin{align}
  & !D_{\fn{diag}}(\gn{!E(x)}, \gn{!A}) \ollabel{repdiag1} \\
  & \lforall[y][(!D_{\fn{diag}}(\gn{!E(x)},y) \lif
  \eq[y][\gn{!A}])]. \ollabel{repdiag2}
\end{align}
Agora mostramos que $\Th{Q} \Proves !A \liff !B(\gn{!A})$. Argumentamos
informalmente, usando apenas a lógica e fatos !!{derivable} em~$\Th{Q}$.

Primeiro, suponha~$!A$, isto é, $!E(\gn{!E(x)})$. Retomando a definição de
$!E(x)$, vemos que $!E(\gn{!E(x)})$ é simplesmente
\[
\lexists[y][(!D_{\fn{diag}}(\gn{!E(x)},y) \land !B(y))].
\]
Considere um tal~$y$. Como $!D_{\fn{diag}}(\gn{!E(x)},y)$, por
\olref{repdiag2}, $y = \gn{!A}$. Logo, de $!B(y)$ obtemos~$!B(\gn{!A})$.

Agora suponha $!B(\gn{!A})$. Por \olref{repdiag1}, temos
\begin{align*}
& !D_{\fn{diag}}(\gn{!E(x)}, \gn{!A}) \land !B(\gn{!A}).
\intertext{Segue-se
que}
& \lexists[y][(!D_{\fn{diag}}(\gn{!E(x)},y) \land !B(y))].
\end{align*}
Mas isso é simplesmente $!E(\gn{!E(x)})$, isto é,~$!A$.
\end{proof}

\begin{digress}
Compare esta prova com a do lema do ponto fixo na teoria da computabilidade.
A diferença é que aqui queremos definir uma \emph{afirmação} em termos de si
mesma, enquanto ali queríamos definir uma \emph{função} em termos de si
mesma; à parte essa diferença, trata-se realmente da mesma ideia.
\end{digress}

\begin{prob} 
  !!^a{formula}~$!A(x)$ é uma \emph{definição de verdade} se $\Th{Q}
  \Proves !B \liff !A(\gn{!B})$ para toda !!{sentence}~$!B$. Mostre, usando
  o lema do ponto fixo, que nenhuma !!{formula} é uma definição de verdade.
\end{prob}

\end{document}

